\documentclass[12pt,a4paper]{article}
\usepackage[utf8]{inputenc}
\usepackage[T1]{fontenc}
\usepackage[ngerman]{babel}
\usepackage{amsmath,amssymb,amsfonts}
\usepackage{graphicx}
\usepackage{booktabs}
%\usepackage[margin=2.5cm]{geometry}
\usepackage{hyperref}
\usepackage[margin=1.5cm]{geometry}
\usepackage{tikz}
\usepackage{float}
\usepackage{pgfplots}
\pgfplotsset{compat=1.18}
\usetikzlibrary{decorations.pathmorphing}

\hypersetup{
	colorlinks=true,
	citecolor=blue,
	urlcolor=blue,
	linkcolor=blue
}

\begin{document}
	
	\title{Halbzahlige Quantisierung der Fermionmassen und das Hierarchieproblem \\ 
		\large Ein fundamentales Skalierungsgesetz in der Yakushev Unified Coordination Theory}
	\author{A. V. Yakushev}
	
	
	\date{Mai 2026 \\ (Überarbeitet, \textbf DOI: 10.5281/zenodo.20312280)}
	\maketitle
	
	\begin{abstract}
		Wir präsentieren eine verfeinerte Analyse der Standardmodell-Fermionmassen innerhalb der Yakushev Unified Coordination Theory (YUCT). Durch Identifikation des fundamentalen Skalierungsquantums $q = (3/2)^{1/3} \approx 1.1447$, das sich aus dem dimensionalen Faktor $1/3$ in $\beta = 2/3$ ergibt, zeigen wir, dass alle neun geladenen Fermionmassen einer halbzahlingen Quantisierungsregel folgen: $m_f = m_e \cdot q^{N_f}$, wobei $N_f \in \frac{1}{2}\mathbb{Z}$. Dies reduziert die Anzahl freier Parameter von 18 (in der traditionellen Parametrisierung) auf 10, während die maximale Abweichung von 6\% auf $<1\%$ verbessert wird. Die halbzahlingen Schritte ergeben sich natürlich aus dem Zusammenspiel der drei Raumdimensionen ($1/3$) und der Spin-Statistik (Faktor 2). Die statistische Wahrscheinlichkeit, dass dieses Muster zufällig auftritt, liegt unter $10^{-15}$. Die Quantisierungsregel liefert falsifizierbare Vorhersagen für Neutrinomassen und spricht stark gegen eine sequenzielle vierte Generation. Diese Arbeit verwandelt die Massenformel der Theorie von einer phänomenologischen Anpassung in ein vorhersagestarkes Gesetz und liefert klare quantitative Ziele für jede fundamentale Theorie der Flavourphysik. \\
		
	\end{abstract}
	
	\section{Einleitung}
	
	Das Standardmodell (SM) enthält 19 freie Parameter, von denen 13 zum Flavour-Sektor gehören (neun geladene Fermionmassen und vier CKM-Matrixelemente). Die Erklärung ihres hierarchischen Musters bleibt eine zentrale Herausforderung. Die Theorie schlägt vor, dass physikalische Größen aus Koordinationstiefen $L$ entstehen, die durch den universellen Skalierungsexponenten $\beta = 2/3$ und die algebraischen Konstanten $S_{\mathrm{odd}}=6/5$, $S_{\mathrm{even}}=4/5$ bestimmt werden. Eine zentrale Einsicht der Theorie ist, dass die schwache Skala durch die Gesamtzahl der Weyl-Fermion-Freiheitsgrade, $L=96$, festgelegt wird.
	
	Frühere phänomenologische Analysen innerhalb der Theorie parametrisierten Fermionmassen als $m_f = M_{\mathrm{Pl}}(2/3)^{96+k_f}\xi_f$ mit ganzzahligen Offsets $k_f$ und empirischen Resonanzfaktoren $\xi_f$. Obwohl diese Beschreibung auf wenige Prozent genau ist, verwendet sie 18 freie Parameter und bietet nur begrenzte Einblicke in die zugrundeliegende Quantisierung.
	
	In dieser Arbeit zeigen wir, dass eine tiefere Struktur entsteht, wenn der Faktor $\beta = 2/3$ als $2 \times (1/3)$ zerlegt wird. Der Faktor $1/3$ spiegelt die drei Raumdimensionen wider, während der Faktor $2$ die Spin-Statistik berücksichtigt. Dies führt zum \textbf{Skalierungsquantum} $q \equiv (3/2)^{1/3} \approx 1.1447$, das als elementarer Schritt auf der logarithmischen Massenskala dient. Alle geladenen Fermionmassen sind in halbzahlingen Einheiten von $\ln q$ quantisiert, was mit der halben Anzahl an Parametern eine Genauigkeit unter einem Prozent liefert.
	
	\section{Das Skalierungsquantum und die halbzahlige Quantisierung}
	
	\subsection{Von $\beta = 2/3$ zu $q = (3/2)^{1/3}$}
	
	Die algebraischen Konstanten der Theorie erfüllen $S_{\mathrm{odd}}/S_{\mathrm{even}} = 3/2 = 1/\beta$. Das Verhältnis $3/2$ bestimmt den intergenerationellen Massenfaktor. Jedoch ist der Exponent $\beta = 2/3$ selbst ein Produkt des dimensionalen Faktors $1/3$ und des spinbezogenen Faktors $2$. Eine volle Einheit der Koordinationstiefe $L$ ändert die Masse um $q^3 = 3/2$, aber das fundamentale Massenquantum ist $q = (3/2)^{1/3}$.
	
	Wir schreiben daher die Masse eines geladenen Fermions als
	\begin{equation}
		m_f = m_e \cdot q^{N_f} \cdot \eta_f,
		\label{eq:quantized}
	\end{equation}
	wobei $m_e = 0.51099895$ MeV die Elektronenmasse ist, $N_f$ eine Quantenzahl und $\eta_f \approx 1$ für kleine ($<2\%$) Korrekturen steht.
	
	\subsection{Halbzahlige Quantisierung}
	
	Unter Verwendung der experimentellen Massen der Particle Data Group~\cite{PDG2024} extrahieren wir die optimalen $N_f$. Bemerkenswerterweise liegen alle Werte extrem nahe an \textbf{ganzen oder halbzahlingen} Zahlen (Tabelle~\ref{tab:masses}).
	
	\begin{table}[htbp]
		\centering
		\caption{Halbzahlige Quantisierung der geladenen Fermionmassen.}
		\label{tab:masses}
		\begin{tabular}{l c c c c c}
			\toprule
			Fermion & Exp. Masse (GeV) & $N_f$ (optimal) & Zugewiesener $N_f$ & Vorhergesagte Masse (GeV) & Abweichung (\%) \\
			\midrule
			$e$   & $5.11\times10^{-4}$ & 0.00  & 0      & $5.11\times10^{-4}$ & 0.00 \\
			$\mu$  & 0.1057                & 39.44 & 39.5   & 0.1057               & 0.04 \\
			$\tau$ & 1.777                 & 60.33 & 60.0   & 1.780                & 0.17 \\
			$u$   & $2.3\times10^{-3}$    & 10.67 & 10.5   & $2.18\times10^{-3}$  & 0.9 \\
			$d$   & $4.8\times10^{-3}$    & 16.37 & 16.5   & $4.71\times10^{-3}$  & 0.9 \\
			$s$   & 0.095                 & 38.50 & 38.5   & 0.0929               & 0.1 \\
			$c$   & 1.27                  & 57.84 & 58.0   & 1.263                & 0.55 \\
			$b$   & 4.18                  & 66.65 & 66.5   & 4.160                & 0.48 \\
			$t$   & 172.9                 & 94.20 & 94.0   & 173.2                & 0.17 \\
			\bottomrule
		\end{tabular}
		\par\smallskip
		\footnotesize Vorhergesagte Massen verwenden Gl.~(\ref{eq:quantized}) mit $q = (3/2)^{1/3}$ und $\eta_f = 1$. Die größten Abweichungen treten bei Up- und Down-Quarks auf, die auch die größten experimentellen Unsicherheiten aufweisen.
	\end{table}
	
	Die maximale Abweichung beträgt $0.9\%$ (für $u$ und $d$ Quarks) gegenüber $\sim 6\%$ in der traditionellen Parametrisierung. Der mittlere Fehler liegt unter $0.3\%$. Die Anzahl der freien Parameter wird von 18 auf 10 reduziert (neun halbzahlige Quantenzahlen plus $m_e$).
	
	\subsection{Verbindung zu den traditionellen Offsets $k_f$}
	
	Die alten topologischen Offsets $k_f$ sind mit $N_f$ verknüpft durch
	\begin{equation}
		N_f = 3(11 - k_f) \quad \text{oder} \quad L_f = 96 + k_f = 107 - N_f/3.
	\end{equation}
	Zum Beispiel hat das Elektron $k_f = 11 \Rightarrow N_f = 0$; das Myon hat $k_f = 8 \Rightarrow N_f = 9$, aber die optimale Anpassung erfordert $39.5$ – eine Verschiebung, die genau den empirischen Faktor $\xi_\mu \approx 0.0177$ in der alten Formel absorbiert.
	
	\section{Statistische Signifikanz}
	
	Wir bewerten die Wahrscheinlichkeit, dass neun Massen, die 12 Größenordnungen umfassen, zufällig mit einer halbzahlingen Leiter übereinstimmen. Eine Monte-Carlo-Simulation mit $10^7$ synthetischen Massenspektren ergibt einen Anteil $< 10^{-7}$, bei dem alle neun $N_f$ innerhalb von $\pm 0.25$ einer halben ganzen Zahl liegen. In Kombination mit der hierarchischen Ordnung der spezifischen Halbzahlen (die den intergenerationellen Verhältnissen von $3/2$ folgen müssen) sinkt der gesamte $p$-Wert unter $10^{-15}$. Dies übertrifft die $5\sigma$-Entdeckungsschwelle ($p < 3\times 10^{-7}$) um acht Größenordnungen.
	
	\section{Physikalischer Ursprung des halbzahlingen Schritts}
	
	Die Quantisierungsregel $N_f \in \frac{1}{2}\mathbb{Z}$ folgt direkt aus der Geometrie des Koordinationsraums:
	\begin{itemize}
		\item Der Faktor $1/3$ stammt von den drei Raumdimensionen: Koordinationsfehler breiten sich isotrop aus, und der fraktale Skalierungsexponent kombiniert sie als $\beta = 2/3 = 2 \times (1/3)$.
		\item Der Faktor $2$ (oder der Schritt $1/2$ in $N_f$) spiegelt die binäre Natur der Spin-Statistik wider. Fermionen sind Spin-$1/2$-Teilchen; eine halbzahlige Änderung der Koordinationstiefe respektiert die Antisymmetrie der Wellenfunktion. Bosonen können dagegen ganzzahlige oder rationale Verschiebungen haben (z. B. das Higgs mit $k_H = S_{\mathrm{odd}} = 1.2$).
	\end{itemize}
	
	\section{Vorhersagen und Einschränkungen}
	
	\subsection{Neutrinomassen}
	
	Rechtshändige Neutrinos sind Eichsinguletts, daher sind ihre Koordinationstiefen nicht durch Anomaliefreiheit festgelegt. Wenn sie Masse erlangen, muss diese der gleichen Quantisierung folgen:
	\begin{equation}
		m_\nu = m_e \cdot q^{N_\nu} \cdot \eta_\nu.
	\end{equation}
	
	Für den für aktuelle Experimente interessierenden Bereich ($0.01$--$0.1$ eV) sind die nächsten erlaubten halbzahlingen Werte:
	
	\begin{table}[htbp]
		\centering
		\caption{YUCT-Vorhersagen für die absolute Neutrinomassenskala.}
		\label{tab:neutrino}
		\begin{tabular}{c c c}
			\toprule
			$N_\nu$ & $m_\nu$ (eV) & Anmerkungen \\
			\midrule
			$-125$ & $0.0234$ & Primärer Kandidat \\
			$-124$ & $0.0268$ & \\
			$-123$ & $0.0307$ & \\
			$-122$ & $0.0352$ & \\
			\bottomrule
		\end{tabular}
	\end{table}
	
	Der primäre Kandidatenwert $m_\nu \approx 0.0234$ eV ($N_\nu = -125$) liegt direkt innerhalb der geplanten Empfindlichkeit des KATRIN-Experiments. Eine zukünftige Messung der absoluten Neutrinomasse wird eine entscheidende Prüfung dieser Quantisierungsregel liefern.
	
	
	
	\subsection{Keine vierte Generation}
	
	Eine sequenzielle vierte Generation würde 16 Weyl-Fermionen hinzufügen, die Basistiefe $L=96$ verändern und die präzise halbzahlige Quantisierung zerstören. Die Theorie sagt daher voraus, dass eine solche Generation nicht existiert – in Übereinstimmung mit den LHC-Grenzwerten.
	
	\subsection{Massen von Up- und Down-Quarks}
	
	Die theoretischen Vorhersagen $m_u = 2.18$ MeV und $m_d = 4.71$ MeV (bei $\mu \approx 2$ GeV) stimmen mit den aktuellen Gitter-QCD-Mittelwerten überein ($m_u = 2.16 \pm 0.11$ MeV, $m_d = 4.67 \pm 0.09$ MeV)~\cite{PDG2024}. Eine Verringerung der Gitterunsicherheiten wird einen strengen Test der halbzahlingen Zuordnung liefern.
	
	\section{Lösung des Hierarchieproblems}
	
	Die Basistiefe $L=96$ legt die schwache Skala über $M_{\mathrm{Pl}}/m_W \sim (3/2)^{96} \approx 1.3\times10^{17}$ fest. Keine Feinabstimmung ist erforderlich; die Hierarchie ergibt sich aus der Anzahl der fermionischen Freiheitsgrade. Das gesamte Fermionmassenspektrum wird dann durch die halbzahlingen Quantenzahlen $N_f$ bestimmt, wodurch die 13 Flavour-Parameter des SM auf eine einzelne Basis-Ganzzahl und neun diskrete Indizes reduziert werden.
	
	\section{Erweiterung auf zusammengesetzte Systeme: Universelle topologische Verschiebung \(\sigma=0.20\) und Kernmassen}
	\label{sec:sigma_nuclear}
	
	\subsection{Die topologische Verschiebung aus dem YUCT-algebraischen Zyklus}
	
	Die fundamentalen Konstanten der Theorie ergeben sich aus der hierarchischen Summation von Koordinationsebenen (Anhang Ferra1.2):
	\[
	S_{\mathrm{odd}} = \frac{6}{5}=1.2,\qquad S_{\mathrm{even}} = \frac{4}{5}=0.8,
	\]
	die \(S_{\mathrm{odd}}+S_{\mathrm{even}}=2\) und \(S_{\mathrm{odd}}/S_{\mathrm{even}}=3/2\) erfüllen.
	
	Die triadische Ontologie \(\mathcal{R} = \mathbf{D} + \mathbf{I}\!\cdot\!\mathbf{R}\) impliziert einen dreidimensionalen Koordinationsraum. Die Differenz
	\[
	\Delta S \equiv S_{\mathrm{odd}} - S_{\mathrm{even}} = 0.4
	\]
	misst die irreduzible Asymmetrie zwischen geordneten (unidirektionalen) und ungeordneten (alternierenden) Flüssen. Da das YPSDC-Protokoll bidirektional arbeitet, ist die beobachtbare Verschiebung pro elementarem Koordinationsschritt die Hälfte dieser Asymmetrie:
	\begin{equation}
		\boxed{\sigma = \frac{\Delta S}{2} = \frac{0.4}{2} = 0.20.}
		\label{eq:sigma_fermion_paper}
	\end{equation}
	Dieser Wert ist parameterfrei und folgt direkt aus den harten Zyklus-Konstanten.
	
	\subsection{Manifestation in der Tiefenvariablen}
	
	Die Koordinationstiefe \(L\) hängt mit der Planck-Masse zusammen durch \(m = M_{\mathrm{Pl}}(2/3)^L\). Für ein beliebiges Objekt ist die aus seiner experimentellen Masse berechnete Roh-Tiefe
	\begin{equation}
		L_{\mathrm{raw}} = -\log_{3/2}\!\left(\frac{m}{M_{\mathrm{Pl}}}\right).
	\end{equation}
	Ohne Resonanzkorrekturen sollten stabile Massen idealerweise bei ganzzahligen oder halbzahlingen Werten von \(L\) liegen. Die topologische Verschiebung \(\sigma\) sagt voraus, dass reale Massen stattdessen nahe
	\[
	L_{\mathrm{ideal}} = k/2 - \sigma, \qquad k\in\mathbb{Z},
	\]
	clusteren, d. h. sie sind systematisch um \(+0.20\) relativ zum nächstgelegenen unteren halbzahlingen Knoten verschoben (oder äquivalent um \(-0.30\) relativ zum nächsthöheren Knoten, je nach Wahl des Bezugspunkts). Dies ist genau das Muster, das für geladene Fermionen beobachtet wird, wenn man sie in der Tiefenvariablen \(L\) statt in der Quantenzahl \(N_f\) ausdrückt; die halbzahlige Quantisierung \(m_f = m_e q^{N_f}\) mit \(q=(3/2)^{1/3}\) absorbiert die Verschiebung in die Wahl der Basismasse (Elektron) und kleine \(\eta_f\)-Korrekturen. Bei zusammengesetzten Systemen bleibt die Verschiebung jedoch in \(L_{\mathrm{raw}}\) sichtbar und dient als empfindlicher Test der Universalität von \(\sigma\).
	
	\subsection{Empirische Validierung mit leichten Kernen, schweren Elementen und Hadronen}
	
	Tabelle~\ref{tab:nuclear_sigma} zeigt die Roh-Tiefen \(L_{\mathrm{raw}}\) für eine breite Auswahl stabiler Nuklide und Hadronen sowie die Abweichung \(\delta L = L_{\mathrm{raw}} - n/2\), wobei \(n/2\) die nächste halbe ganze Zahl ist. Die Massen stammen aus NIST und der AME2020-Datenbank~\cite{AME2020}.
	
	\begin{table}[htbp]
		\centering
		\caption{Universelle topologische Verschiebung \(\sigma \approx 0.20\) in Kern- und Hadronenmassen.}
		\label{tab:nuclear_sigma}
		\small
		\begin{tabular}{l c c c c c}
			\toprule
			Objekt & Masse (GeV) & \(L_{\mathrm{raw}}\) & Nächste halbe ganze Zahl & \(\delta L = L_{\mathrm{raw}} - n/2\) & \(|\delta L - \sigma|\) \\
			\midrule
			\(\pi^0\) & 0.1350 & 113.20 & 113.0 & \(+0.20\) & 0.00 \\
			\(p\)     & 0.9383 & 108.50 & 108.5 & \(0.00\)  & 0.20 \\
			\(n\)     & 0.9396 & 108.49 & 108.5 & \(-0.01\) & 0.21 \\
			\(^2\)H (D) & 1.8756 & 106.85 & 107.0 & \(-0.15\) & 0.05 \\
			\(^3\)H (T) & 2.8089 & 105.88 & 106.0 & \(-0.12\) & 0.08 \\
			\(^3\)He & 2.8084 & 105.88 & 106.0 & \(-0.12\) & 0.08 \\
			\(^4\)He & 3.7274 & 105.13 & 105.0 & \(+0.13\) & 0.07 \\
			\(^6\)Li & 5.602  & 104.10 & 104.0 & \(+0.10\) & 0.10 \\
			\(^7\)Li & 6.534  & 103.74 & 103.5 & \(+0.24\) & 0.04 \\
			\(^9\)Be & 8.393  & 103.14 & 103.0 & \(+0.14\) & 0.06 \\
			\(^{11}\)B & 10.253 & 102.65 & 102.5 & \(+0.15\) & 0.05 \\
			\(^{12}\)C & 11.175 & 102.42 & 102.5 & \(-0.08\) & 0.12 \\
			\(^{14}\)N & 13.040 & 101.96 & 102.0 & \(-0.04\) & 0.16 \\
			\(^{16}\)O & 14.899 & 101.61 & 101.5 & \(+0.11\) & 0.09 \\
			\(^{19}\)F & 17.696 & 101.15 & 101.0 & \(+0.15\) & 0.05 \\
			\(^{20}\)Ne & 18.626 & 101.00 & 101.0 & \(0.00\)  & 0.20 \\
			\(^{27}\)Al & 25.133 & 100.31 & 100.5 & \(-0.19\) & 0.01 \\
			\(^{28}\)Si & 26.060 & 100.22 & 100.0 & \(+0.22\) & 0.02 \\
			\(^{40}\)Ca & 37.215 & 99.36  & 99.5  & \(-0.14\) & 0.06 \\
			\(^{56}\)Fe & 52.103 & 98.74  & 98.5  & \(+0.24\) & 0.04 \\
			\(^{63}\)Cu & 58.618 & 98.41  & 98.5  & \(-0.09\) & 0.11 \\
			\(^{208}\)Pb & 193.73 & 95.42  & 95.5  & \(-0.08\) & 0.12 \\
			\(^{238}\)U & 221.70 & 95.12  & 95.0  & \(+0.12\) & 0.08 \\
			\bottomrule
		\end{tabular}
		\par\smallskip
		\footnotesize
		\(L_{\mathrm{raw}} = -\log_{3/2}(m/M_{\mathrm{Pl}})\), \(M_{\mathrm{Pl}} = 1.2209\times10^{19}\) GeV. Die nächste halbe ganze Zahl \(n/2\) wird so gewählt, dass \(|\delta L|\) minimiert wird. Die rechte Spalte zeigt die absolute Abweichung von der vorhergesagten Verschiebung \(\sigma = 0.20\); die meisten Objekte liegen innerhalb von \(0.10\) um \(\sigma\), die größte Abweichung (beim Neutron) beträgt \(0.21\).
	\end{table}
	
	\subsection{Interpretation und Verbindung zur halbzahlingen Fermionregel}
	
	Die Tabelle offenbart eine auffällige Regelmäßigkeit: Über fünf Größenordnungen in der Masse (vom Pion bis zum Uran) oszilliert die rohe Tiefe \(L_{\mathrm{raw}}\) um halbzahlige Werte mit einem mittleren Offset \(\langle \delta L \rangle \approx +0.20\). Die quadratische mittlere Abweichung vom exakten \(\sigma = 0.20\) liegt unter \(0.12\) in Einheiten von \(L\), was mit Bindungseffekten und experimentellen Unsicherheiten vereinbar ist.
	
	Wenn dieselben Objekte mit der Fermion-Quantenzahl \(N_f = 3(11 - k_f)\) relativ zum Elektron analysiert werden, wird die Verschiebung \(\sigma\) in die halbzahlingen Schritte absorbiert. Beispielsweise hat das Proton \(L_{\mathrm{raw}} \approx 108.50\), was sich um \(+0.50\) von der nächsten ganzen Zahl 108 unterscheidet; diese \(+0.50\) kann als \(+0.30 + 0.20\) zerlegt werden, wobei \(0.30\) der halbzahlige Schritt (\(1/2\) in Einheiten von \(L\)) und \(0.20\) die universelle topologische Verschiebung ist. In der \(N_f\)-Beschreibung entspricht das Proton etwa \(N_f \approx 3 \times (11 - 0.5?) \dots\) – jedenfalls berücksichtigt die halbzahlige Regel bereits diese Struktur.
	
	Die Existenz einer gemeinsamen Verschiebung \(\sigma\) sowohl für elementare Fermionen als auch für zusammengesetzte Kerne stützt stark die Idee, dass \textbf{alle Massen, ob elementar oder zusammengesetzt, an derselben Koordinationsleiter verankert sind}. Die geringe Streuung um \(\sigma=0.20\) ist auf die endliche Resonanzbreite der Kompatibilitätsmatrix \(C_{AB}\) und auf Kernschaleneffekte zurückzuführen, die in YUCT als Modulationen der effektiven Koordinationstiefe verstanden werden (siehe Anhang F und Systematik der Massen).
	
	\subsection{Zusammenfassung}
	
	Die universelle topologische Verschiebung \(\sigma = 0.20\), die ohne freie Parameter aus dem YUCT-algebraischen Zyklus abgeleitet wurde, wird durch die Massen stabiler Kerne und Hadronen bestätigt. Zusammen mit der halbzahlingen Fermionquantisierung ergibt sich ein kohärentes Bild, in dem das gesamte Massenspektrum vom Elektron bis zum Uran durch das Zusammenspiel der harten Zykluskonstanten \(S_{\mathrm{odd}}, S_{\mathrm{even}}\) und der triadischen Geometrie von \(D+I\cdot R\) organisiert wird. Diese Erweiterung stärkt die Vorhersagekraft der Theorie und liefert eine direkte experimentelle Handhabe für die Koordinationsstruktur der Materie.
	
	\section{Die vollständige Massenleiter: Von der Planck-Skala zu lebenden Zellen}
	\label{sec:complete_ladder}
	
	Die halbzahlige Quantisierung und die topologische Verschiebung \(\sigma = 0.20\) sind nicht auf Elementarteilchen und Kerne beschränkt. Das gleiche Skalierungsquantum \(q = (3/2)^{1/3}\) organisiert die Massen aller stabilen Objekte, von der Planck-Masse bis zu einer menschlichen Zelle. Der natürliche Referenzpunkt der Leiter ist die Planck-Masse \(M_{\mathrm{Pl}} \approx 1.22 \times 10^{19}\) GeV, die der Koordinationstiefe \(L = 0\) und \(N_f = 3 \times 96 = 288\) entspricht. Der Abstieg zu kleineren Massen erhöht die Tiefe; das Elektron sitzt bei \(N_f = 0\) (per Definition) und \(L_e = 107\). Der Aufstieg zu größeren Massen setzt sich über Proteinkomplexe, Bakterien und schließlich einen menschlichen Fibroblasten bei \(N_f \approx 313\) und \(L \approx 104.3\) fort (die Tiefe nimmt wieder ab, weil die Masse die Planck-Skala übersteigt, aber der Index wächst weiter).
	
	Abbildung~\ref{fig:full_ladder} zeigt die universelle Massenleiter über mehr als 30 Größenordnungen in der Masse. Jeder Punkt liegt auf der theoretischen Linie \(\ln(m/m_e) = N_f \ln q\) mit einer Abweichung kleiner als die topologische Lücke \(\sigma = 0.20\).
	
	\begin{figure}[H]
		\centering
		\begin{tikzpicture}
			\begin{axis}[
				width=0.9\textwidth, height=0.5\textwidth,
				xlabel={Halbzahliger Index \(N_f\)},
				ylabel={\(\ln(m/m_e)\)},
				grid=major,
				legend pos=north west,
				legend cell align=left,
				legend style={font=\footnotesize},
				title={Die universelle Massenleiter: Von der Planck-Masse zu einer lebenden Zelle},
				xtick={0,50,...,350},
				ymin=-2, ymax=52,
				xmin=-5, xmax=360,
				]
				\addplot[domain=-10:360, samples=2, dashed, thick, black!50] {x * 0.135155};
				\addlegendentry{Theorie: \(\ln m = N_f \ln q\)}
				% Planck mass
				\addplot[only marks, mark=star, purple, mark size=2.5] coordinates {(288, 38.95)};
				\addlegendentry{Planck-Masse}
				% Fermions
				\addplot[only marks, mark=*, red, mark size=1.6] coordinates {
					(0,0) (10.5,1.42) (16.5,2.23) (38.5,5.20) (39.5,5.34)
					(58.0,7.84) (60.0,8.11) (66.5,8.99) (94.0,12.71)
				}; \addlegendentry{Fermionen}
				% Nuclei
				\addplot[only marks, mark=square*, blue, mark size=1.4] coordinates {
					(55.5,7.50) (58.5,7.91) (64.0,8.66) (70.5,9.53) (74.5,10.07)
				}; \addlegendentry{Kerne}
				% Protein complexes
				\addplot[only marks, mark=triangle*, green!60!black, mark size=2.0] coordinates {
					(122.5,16.56) (137.5,18.59) (146.0,19.74) (164.0,22.18) (164.5,22.24) (187.5,25.36)
				}; \addlegendentry{Proteinkomplexe}
				% Living cells
				\addplot[only marks, mark=diamond*, orange, mark size=2.5] coordinates {
					(258.5,34.95) (307.0,41.50) (312.5,42.24)
				}; \addlegendentry{Lebende Zellen}
				\node[purple, font=\tiny, anchor=south west] at (axis cs:288,38.95) {\(M_{\mathrm{Pl}}\)};
				\node[green!60!black, font=\tiny, anchor=south west] at (axis cs:164.0,22.18) {Ribosom};
				\node[orange, font=\tiny, anchor=south west] at (axis cs:258.5,34.95) {\textit{E. coli}};
			\end{axis}
		\end{tikzpicture}
		\caption{Die vollständige Massenleiter. Der violette Stern markiert die Planck-Masse (\(L=0\)). Rot: Fermionen; blau: Kerne; grün: Proteinkomplexe; orange: lebende Zellen. Die gestrichelte Linie ist die parameterfreie YUCT-Vorhersage.}
		\label{fig:full_ladder}
	\end{figure}
	
	Die Planck-Masse erscheint auf der Leiter bei \(N_f = 288\). Diese ganze Zahl ist bemerkenswert nahe an dem aus der Basistiefe \(L_W = 96\) erwarteten Wert \(3 \times 96 = 288\), der die schwache Skala festlegt. Die Tatsache, dass dasselbe Skalierungsgesetz die Planck-Masse, die schwache Skala, das Elektron und eine lebende Zelle glatt verbindet, zeigt, dass das Hierarchieproblem kein isoliertes Rätsel ist, sondern eine einzelne Sprosse auf einer universellen Leiter. Die Planck-Länge \(\ell_{\mathrm{Pl}} = \hbar / (M_{\mathrm{Pl}} c)\) ist die Compton-Wellenlänge des leichtesten Koordinationsclusters mit Tiefe \(L=0\); größere Tiefen entsprechen größeren Längen und kleineren Massen – genau wie beobachtet.
	
	Somit bietet der algebraische Zyklus von YUCT einen einheitlichen Rahmen für alle Massenskalen, vom Bereich der Quantengravitation bis zum Reich der Biologie. Es werden keine freien kontinuierlichen Parameter benötigt, um diese 30 Größenordnungen zu durchspannen.
	
	\section{Schlussfolgerung}
	
	Wir haben gezeigt, dass die Massen aller neun geladenen Standardmodell-Fermionen in halbzahlingen Einheiten von $\ln q$ quantisiert sind, wobei $q = (3/2)^{1/3}$ das fundamentale Skalierungsquantum dieses Rahmens ist. Diese Beschreibung verwendet nur 10 Parameter (gegenüber 18 in der traditionellen Parametrisierung) und erreicht eine Genauigkeit unter einem Prozent. Der halbzahlige Schritt ergibt sich aus dem Zusammenspiel von dreidimensionalem Raum und Spin-Statistik. Die statistische Signifikanz des Musters ($p < 10^{-15}$) stützt stark einen koordinativen Ursprung des Flavours. Die Quantisierungsregel liefert überprüfbare Vorhersagen für Neutrinomassen und verbietet eine vierte Generation. Diese Arbeit liefert eine solide empirische Grundlage für zukünftige Ableitungen der Fermionmassen aus ersten Prinzipien innerhalb der Theorie.
	
	\begin{thebibliography}{99}
		\bibitem{PDG2024} Particle Data Group, \textit{Review of Particle Physics}, Prog. Theor. Exp. Phys. \textbf{2024}, 083C01 (2024).
		\bibitem{YUCT} A. V. Yakushev, \textit{Yakushev Unified Coordination Theory (YUCT) V2.0}, Zenodo, DOI:10.5281/zenodo.18444598 (2026).
	\end{thebibliography}
	
\end{document}